\section{Методы и организация исследования}

\subsection{Подтверждённая архитектура исходного прототипа}

Предыдущая НИР реализована с помощью PyTorch, Hugging Face Transformers, Ray Actors, FastAPI и Docker Compose. В качестве небольшой контрольной модели используется \code{sshleifer/tiny-gpt2}: $N=2$ Transformer-блока, $H=2$, $V=50\,257$. В рамках одного физического компьютера запускаются контейнеры \code{ray-head}, \code{ray-worker-1}, \code{ray-worker-2}, зарегистрированные как три логических Ray-узла.

Исполнитель 0 содержит эмбеддинги и блок $[0,1)$, исполнитель 1~--- блок $[1,2)$, исполнитель 2~--- заключительную нормализацию и \code{lm\_head} при пустом диапазоне $[2,2)$. На каждом шаге генерации через удалённые вызовы последовательно проходят промежуточные тензоры. В исходном варианте отсутствует KV-кэш: текущая последовательность обрабатывается целиком на каждом шаге. Это упрощает контроль корректности, хотя и увеличивает вычислительные затраты при росте длины контекста.

\begin{figure*}[t]
\centering
\begin{tikzpicture}[every node/.style={font=\footnotesize},node distance=5mm,
box/.style={draw,rounded corners,align=center,minimum height=9mm,text width=5.0cm,fill=gray!6},
arr/.style={-{Stealth[length=2mm]},line width=0.65pt}]
\node[box] (a) {Ray head / shard 0\\эмбеддинги + Transformer-блок 0};
\node[box,below=of a] (b) {Ray worker 1 / shard 1\\Transformer-блок 1};
\node[box,below=of b] (c) {Ray worker 2 / shard 2\\финальная нормализация + lm\_head};
\node[draw,rounded corners,align=center,text width=3.0cm,right=18mm of b,minimum height=16mm] (d) {Управляющий процесс\\выбор / добавление токена};
\draw[arr] (a)--node[right]{hidden states}(b);
\draw[arr] (b)--node[right]{hidden states}(c);
\draw[arr] (c.east)-|node[pos=0.24,right,font=\scriptsize]{logits / token\_id}(d.south);
\draw[arr,dashed] (d.north)|- (a.east);
\node[font=\scriptsize,align=center,below=6mm of c] {Три контейнера Ray на \textbf{одном} физическом компьютере};
\end{tikzpicture}
\caption{Исполнение шагов генерации в исходной архитектуре; режимы D1--D3 изменяют только обработку выходных логитов.}
\label{fig:pipeline}
\end{figure*}

По материалам выполненной НИР на таком стенде подтверждены три активных Ray-узла и успешная генерация 20 токенов. При этом 24 CPU, которые отображались как суммарные ресурсы Ray, являются \emph{логическими объявлениями ресурсов}, а не 24 независимыми физическими ядрами \cite{rayresources}. При переносе на MacBook Pro с Apple M1 Pro и 16~ГБ памяти параметры Docker и фактическая загрузка должны быть проверены повторно: прежние значения скорости нельзя автоматически распространять на новый компьютер.

\subsection{Модель коммуникационных расходов}

Пусть $B$~--- размер пакета, $d_z$~--- число байт на один логит, $T$~--- число генерируемых токенов, $L_t=L_0+t-1$~--- длина контекста на шаге $t$. Для режима D0 сырой объём \emph{полезных выходных данных} от завершающего исполнителя составляет
\begin{equation}
 V_{\mathrm{D0}} = B V d_z\sum_{t=1}^{T}L_t
 = BVd_z\left(TL_0+\frac{T(T-1)}{2}\right).
 \label{eq:full}
\end{equation}
Для D1 и D2 выходной объём одинаков, хотя стоимость вычислений различается:
\begin{equation}
 V_{\mathrm{D1}}=V_{\mathrm{D2}}=B V d_z T.
 \label{eq:last}
\end{equation}
Для D3 выход состоит из идентификаторов токенов длиной $d_{\mathrm{id}}$ байт:
\begin{equation}
 V_{\mathrm{D3}}=B d_{\mathrm{id}} T.
 \label{eq:id}
\end{equation}

Для выполненного ранее контрольного сценария $B=1$, $L_0=5$, $T=20$, $V=50\,257$ и $d_z=4$ получаем $\sum L_t=290$. Тогда
\begin{equation}
 V_{\mathrm{D0}}=58\,298\,120\ \text{байт}\approx55{,}60\ \text{МиБ},
 \label{eq:example}
\end{equation}
а для D1--D2~--- $4\,020\,560$ байт (примерно $3{,}83$~МиБ). При передаче идентификаторов \code{int64} в D3~--- $160$ байт. Все три величины относятся \emph{только к выходному направлению} и являются расчётными, а не замерами транспортного трафика.

Для промежуточных hidden states на одной границе при последовательной передаче полного контекста величина составляет $Bd_hH\sum_t L_t$, где $d_h$~--- размер элемента активации. В исходной НИР суммарный объём двух границ оценён в $4\,640$ байт, а отдельно передаваемых масок~--- ещё в $4\,640$ байт. Поэтому минимальный сырой объём всех рассматриваемых тензоров в D0 равен $58\,307\,400$ байт (около $55{,}61$ МиБ) без заголовков Ray, сериализации и дополнительных сообщений.

Модель задержки должна учитывать как минимум
\begin{equation}
 t_{\mathrm{gen}}\approx t_{\mathrm{compute}}+t_{\mathrm{rpc}}
 +t_{\mathrm{serialize}}+t_{\mathrm{transport}}+t_{\mathrm{schedule}}.
 \label{eq:latency}
\end{equation}
Эта декомпозиция является \emph{моделью для измерений}, а не установленным экспериментально равенством: компоненты могут перекрываться. Уменьшение $V_{\mathrm{payload}}$ не гарантирует пропорционального уменьшения $t_{\mathrm{gen}}$.

\subsection{Разделение вычислительного и коммуникационного эффекта}

Предлагается последовательная абляция D0--D3. При переходе D0$\to$D1 сокращается объём возвращаемого сообщения без обязательного сокращения вычисления выходной головы. При D1$\to$D2 логиты вычисляются только для последней позиции: для причинной GPT-2-модели это допускается послойной нормализацией и независимым применением выходной линейной головы к каждой позиции. При D2$\to$D3 выбор \code{argmax} переносится внутрь последнего Ray Actor, и вместо вектора возвращается один ID.

Каждый переход проверяется на идентичность последовательностей токенов эталонному B0 при зафиксированных весах и численной конфигурации. Вклад отдельных изменений оценивается попарными сравнениями D0--D1, D1--D2, D2--D3. Такое разложение позволяет не приписывать экономии передачи данных ускорение, фактически достигнутое сокращением вычислений.

\subsection{Математическая задача выбора послойного разбиения}

Для дальнейшей работы с многослойной GPT-2-совместимой моделью вводится отдельная оптимизационная постановка. Пусть $N\ge K$, блоки сохраняют исходный порядок, а $\omega$ обозначает один нагрузочный сценарий (длину входа, генерации и размер пакета). Для исполнителя $i$ определяются оценённое время вычислений $C_i(p,q;\omega)$ для блоков $[p,q)$, потребность в памяти $M_i(p,q;\omega)$ и доступный предел $M_i^{\max}$. Эмбеддинги относятся к первому, а выходная голова~--- к последнему исполнителю и учитываются в их функциях стоимости.

Для последовательного исполнения вводится аддитивная \emph{прогнозная} стоимость
\begin{equation}
 \widehat T(\mathbf b;\omega)
 =\sum_{i=1}^{K} C_i(b_{i-1},b_i;\omega)
 +\sum_{i=1}^{K-1} Q_i(b_i;\omega),
 \label{eq:cost}
\end{equation}
где $Q_i$~--- стоимость передачи активаций на границе между исполнителями, в том числе задержка вызова, сериализация и оценка передачи данных. Для \emph{независимых} физических каналов возможна приближённая форма $Q_i\approx\alpha_i+S_i/\beta_i$, где $S_i$~--- число байт, $\beta_i$~--- эффективная скорость линии; для локальных контейнеров параметры следует измерять как межпроцессные затраты, а не выдавать за Ethernet-характеристики.

На конечном наборе нагрузок $\Omega$ с весами $w_\omega\ge0$ и $\sum_\omega w_\omega=1$ оптимизационная задача имеет вид
\begin{align}
 \min_{0=b_0<b_1<\cdots<b_K=N}\;&\sum_{\omega\in\Omega}w_\omega\widehat T(\mathbf b;\omega),\label{eq:opt}\\[-2mm]
 \text{при условиях}\quad &M_i(b_{i-1},b_i;\omega)\le M_i^{\max},\quad \forall i,\omega.\notag
\end{align}
Таким образом, оптимизация направлена на среднюю прогнозную \emph{задержку одного запроса} по заданной нагрузке, а не на несуществующее в исходном прототипе перекрытие микропакетов. Если ограничения памяти нарушаются, соответствующее разбиение недопустимо.

\subsection{Метод решения и проверка оптимальности в модели}

Поскольку порядок блоков фиксирован и стоимость~(\ref{eq:cost}) аддитивна по стадиям и границам, для модели предлагается динамическое программирование. Обозначим через $F(i,q)$ минимальную допустимую среднюю стоимость размещения первых $q$ блоков на первых $i$ исполнителях. Для каждой допустимой границы $p<q$:
\begin{equation}
\begin{split}
 F(i,q)=\min_{p}\bigl[&F(i-1,p)+\overline C_i(p,q)\\
 &+\overline Q_{i-1}(p)\bigr],
\end{split}
 \label{eq:dp}
\end{equation}
где $\overline C_i=\sum_\omega w_\omega C_i(\cdot;\omega)$ и аналогично $\overline Q_{i-1}$, а варианты с $M_i(p,q;\omega)>M_i^{\max}$ исключаются. Для первого исполнителя $Q_0=0$; $F(0,0)=0$, остальные недопустимые начальные состояния бесконечны.

\begin{proposition}
При фиксированном порядке узлов, аддитивной стоимости~(\ref{eq:cost}) и проверяемых для каждого сегмента ограничениях памяти рекуррентное соотношение~(\ref{eq:dp}) находит минимальную стоимость среди всех допустимых непрерывных разбиений модели.
\end{proposition}
\noindent\emph{Обоснование.} В любом оптимальном разбиении последняя стадия покрывает некоторый отрезок $[p,q)$. Если размещение предшествующих $p$ блоков не оптимально для первых $i-1$ исполнителей, его можно заменить более дешёвым допустимым размещением, уменьшив общую аддитивную стоимость. Поэтому оптимум обязан содержать оптимальную подзадачу и учитывается перебором всех $p$. По индукции получаем рекуррентное соотношение.

При заранее подготовленных стоимостях сегментов вычислительная сложность такого алгоритма составляет $O(KN^2)$, память для таблицы состояний~--- $O(KN)$. Для небольших $N$ корректность результата планируется проверять независимым полным перебором границ: например, при $N=6$, $K=3$ существует $\binom{5}{2}=10$ способов разделить блоки на три непустые последовательные группы.

Оптимальность полученного решения гарантируется только \emph{в рамках заданной модели затрат}. На MacBook, где три контейнера конкурируют за общие физические ядра и память, суммарное время может не быть аддитивным. Поэтому прогноз модели необходимо независимо сравнивать с фактическими замерами, а при сильном расхождении~--- пересматривать профилирование и ограничения. Для многозапросного перекрывающегося конвейера целевая функция будет другой; такая функциональность относится к дальнейшему развитию, а не к выполненному прототипу \cite{hwang2026}.

\subsection{Проект авторского набора данных и проверок}

Предусматривается собственный воспроизводимый набор из 144 нагрузочных сценариев: четыре категории текстов (продолжение фразы, вопросы, техническое описание, инструкции), три исходных текста в категории, четыре целевые длины входа (16, 64, 128, 256 токенов) и три длины генерации (8, 16, 32 токена). Число конфигураций равно $4\cdot3\cdot4\cdot3=144$. Тексты и правила их подготовки будут сохранены в файлах данных и генераторе; фактические длины определяются GPT-2-токенизатором, а не приписываются исходным текстам по целевому значению. Для каждой записи планируется фиксировать \code{case\_id}, \code{prompt\_text}, \code{input\_tokens}, \code{max\_new\_tokens}, версию модели и происхождение текста.

Набор не предназначен для обучения модели: это \emph{авторские входные нагрузки для измерения свойств системы}. Отдельно предполагается сохранять машиночитаемый журнал измерений по всем режимам, в том числе ошибочные запуски и несовпадения токенов. Исследовательский датасет и результаты измерений на момент ПЗ-3 ещё не созданы; здесь представлена проверяемая спецификация их получения.

\subsection{Экспериментальная методика в пределах доступных ресурсов}

Планируемая базовая аппаратная платформа~--- MacBook Pro с Apple M1 Pro, 16~ГБ объединённой памяти и приблизительно 146~ГБ свободного пространства. Docker Desktop на Mac использует Linux-виртуальную машину с отдельным лимитом оперативной памяти, по умолчанию 50\% памяти хоста \cite{dockermemory}. Следовательно, реальная доступность ресурсов для Ray определяется настройками Docker и нагрузкой macOS. В Linux-контейнерах нельзя автоматически рассчитывать на ускорение PyTorch через macOS MPS; основная серия задумывается как CPU-эксперимент.

Развитие модели выполняется ступенчато: сначала воспроизводятся базовые тесты на tiny-gpt2, далее проверяется \code{distilbert/distilgpt2} (приблизительно 82 млн параметров по карточке модели), затем при достаточном запасе памяти~--- \code{openai-community/gpt2} (124 млн параметров) \cite{distilgpt2card,gpt2card}. При переходе необходимо избежать пикового дублирования полных весов, которое может возникнуть, если каждый Ray Actor сначала загружает всю модель, а затем удаляет ненужные части.

Для каждого этапа применяется последовательность проверок: \emph{(1)} запуск и проверка памяти на небольшой пилотной выборке; \emph{(2)} точная проверка идентификаторов токенов и диапазонов слоёв; \emph{(3)} повторные измерения после прогрева; \emph{(4)} расширение нагрузки только при успешных проверках. Времена сравниваются на одной физической машине, с фиксированными версиями библиотек, числом потоков PyTorch и ограничениями Ray Actors. Журналы должны содержать фактические значения CPU, RAM, версию окружения и ошибки. Запуски с нехваткой памяти не исключаются незаметно из анализа.

\begin{table*}[t]
\centering\small
\caption{Этапность и критерии допустимости исследований}
\label{tab:plan}
\begin{tabularx}{\textwidth}{@{}p{3.2cm}X p{3.0cm}@{}}
\toprule
\textbf{Этап} & \textbf{Критерий перехода} & \textbf{Статус}\\
\midrule
Исходный tiny-gpt2 / Ray & Успешная генерация, размещение трёх частей, базовые тесты & Реализовано в НИР\\
Четыре варианта логитов & Согласованные token ID, учёт вычислений и обмена & Планируется\\
Авторские 144 сценария & Проверенные токены, происхождение, версии и журналы & Планируется\\
DistilGPT2, затем GPT-2 & Устойчивый CPU-запуск без переполнения памяти & Условно планируется\\
Оптимизация разбиения & Сверка DP с перебором и независимым замером & Планируется\\
KV-кэш, многозапросный pipeline & Отдельная архитектурная и экспериментальная работа & Перспектива\\
Несколько физических серверов & Наличие дополнительного оборудования и сетевых замеров & Вне текущего стенда\\
\bottomrule
\end{tabularx}
\end{table*}

На стадии сравнения необходимо использовать одинаковые сценарии для всех режимов, фиксировать прогрев, повторять запуски и оценивать медиану и разброс задержек. Измерения полезного объёма тензоров следует отделять от инструментального измерения сетевых байтов и сериализации. Для моделей с одинаковой математикой допустимо ожидать совпадения токенов в детерминированном режиме, но факт совпадения должен быть подтверждён проверкой, а не объявлен заранее.

\subsection{Воспроизводимость и границы применимости}

Результирующий экспериментальный проект должен включать исходники, фиксированные зависимости, конфигурации запуска, параметры аппаратной среды, генератор нагрузочного датасета, идентификаторы моделей, сохранённые метрики и скрипты построения графиков. Теоретические оценки, прогнозы модели распределения и измеренные показатели должны маркироваться раздельно.

Следует различать: (а) доказанное корректное выполнение учебного распределённого прототипа; (б) математическую возможность сократить передаваемое сообщение; (в) будущую проверку фактического ускорения; (г) дальнейшую задачу оптимального распределения слоёв. Уменьшение объёма данных не равно установленному ускорению, а оптимальность в аддитивной модели не равна оптимальности на реальном оборудовании. С учётом этих ограничений предлагаемый план остаётся реализуемым на доступной рабочей станции без обязательного использования GPU-кластера.
